%==============================================================================
\section{Thuật toán Random Forest}
\label{sec:random-forest}
%==============================================================================

\begin{dinhnghia}[title={Random Forest}]
Thuật toán học máy dùng \textbf{tập hợp} (ensemble) cây quyết định để dự đoán; Leo Breiman, 2001.
Tạo nhiều cây, mỗi cây huấn luyện trên tập con dữ liệu khác nhau; kết hợp dự đoán.
\end{dinhnghia}

\subsection{Cách Random Forest hoạt động}

\begin{phuongphap}[title={Các bước}]
\begin{enumerate}
  \item Chọn ngẫu nhiên mẫu từ dataset.
  \item Xây cây quyết định cho từng mẫu; thu dự đoán từng cây.
  \item \textbf{Bỏ phiếu} (voting) cho từng kết quả dự đoán.
  \item Chọn kết quả được bỏ phiếu nhiều nhất làm dự đoán cuối.
\end{enumerate}
\end{phuongphap}

\tsimg{07_random_forest/img01.jpg}

\begin{ynghia}[title={Phân loại và hồi quy}]
Phân loại: dùng \textbf{mode} (lớp đa số) của các cây.
Hồi quy: dùng \textbf{trung bình} dự đoán các cây.
\end{ynghia}

\subsection{Ưu điểm}

\begin{tinhchat}[title={Ưu điểm chính}]
\begin{itemize}
  \item \textbf{Chống overfitting}: ensemble giảm ảnh hưởng outlier và nhiễu.
  \item \textbf{Độ chính xác cao}: kết hợp nhiều cây.
  \item \textbf{Xử lý dữ liệu thiếu}: không bắt buộc mọi feature có mặt ở mọi điểm.
  \item \textbf{Quan hệ phi tuyến}: cây quyết định mô hình hóa phi tuyến.
  \item \textbf{Feature importance}: hỗ trợ chọn/làm feature.
\end{itemize}
\end{tinhchat}

\subsection{Triển khai Python}

\begin{vidu}[title={Bước 1--2: Import và nạp Iris}]
\begin{lstlisting}
import pandas as pd
from sklearn.ensemble import RandomForestClassifier
from sklearn.datasets import load_iris

iris = load_iris()
X = pd.DataFrame(iris.data, columns=iris.feature_names)
y = iris.target
\end{lstlisting}
\end{vidu}

\begin{vidu}[title={Bước 3: Tiền xử lý và chia dữ liệu}]
\begin{lstlisting}
from sklearn.model_selection import train_test_split
X_train, X_test, y_train, y_test = train_test_split(
    X, y, test_size=0.35, random_state=42)
\end{lstlisting}
\end{vidu}

\begin{vidu}[title={Bước 4--5: Huấn luyện và dự đoán}]
\begin{lstlisting}
rfc = RandomForestClassifier(n_estimators=100)
rfc.fit(X_train, y_train)
y_pred = rfc.predict(X_test)
\end{lstlisting}
\end{vidu}

\begin{vidu}[title={Bước 6: Đánh giá}]
\begin{lstlisting}
from sklearn.metrics import (accuracy_score, precision_score,
                             recall_score, f1_score)
print("Accuracy:", accuracy_score(y_test, y_pred))
print("Precision:", precision_score(y_test, y_pred, average='weighted'))
print("Recall:", recall_score(y_test, y_pred, average='weighted'))
print("F1-score:", f1_score(y_test, y_pred, average='weighted'))
\end{lstlisting}

\textbf{Output mẫu:}
\begin{verbatim}
Accuracy: 0.9811320754716981
Precision: 0.9821802935010483
Recall: 0.9811320754716981
F1-score: 0.9811157396063056
\end{verbatim}
\end{vidu}

\subsection{Ưu và nhược điểm}

\begin{tinhchat}[title={Pros}]
Giảm overfitting nhờ trung bình/kết hợp cây; hoạt động tốt trên nhiều loại dữ liệu hơn một cây; phương sai thấp hơn một cây; linh hoạt, chính xác cao; \textbf{không bắt buộc scale} feature.
\end{tinhchat}

\begin{luuy}[title={Cons}]
Phức tạp cao; xây dựng tốn thời gian và tài nguyên hơn một cây; khó trực quan khi có rất nhiều cây; dự đoán chậm hơn một số thuật toán khác.
\end{luuy}
